% !PRINTED 138
\weberpage{138}

\begin{center}
{\Large Fourth Section.}\\[0.8em]
{\Large\gesperrt{Symmetric Functions.}}\\[0.8em]
\rule{0.16\textwidth}{0.4pt}
\end{center}

\begin{center}
{\large \S.\ 41.}

\vspace{0.5em}
{\large\gesperrt{Concept of Symmetric Functions.}}

{\large\gesperrt{Symmetric Fundamental Functions.}}
\end{center}

In this section we consider \gesperrt{whole functions}, of arbitrary degree, in an arbitrary number $n$ of variables
\[
 \alpha_1,\alpha_2,\ldots,\alpha_n.
\]

Such a function is called \gesperrt{symmetric} \gesperrt{if it remains unchanged when the variables $\alpha_1,\alpha_2,\ldots,\alpha_n$ are subjected to an arbitrary permutation;} and it is such symmetric functions whose properties and laws of formation we now have to learn more precisely.

In order that a function $\Phi(\alpha_1,\alpha_2,\ldots,\alpha_n)$ be symmetric, it is sufficient that it not be changed when any two of the arguments $\alpha_1,\alpha_2,\ldots,\alpha_n$ are interchanged, since by \S\,18 all permutations can be formed by a succession of transpositions.

The function $\Phi$ will in general not be homogeneous, but will contain terms of different dimensions. If, however, one collects all terms of the same dimension, every $\Phi$ can be represented as a sum of \gesperrt{homogeneous functions} of different degrees; and, if $\Phi$ is to be symmetric, then each homogeneous constituent of a definite degree must by itself be symmetric, since the dimensions of the terms are not changed by the permutations of the variables.
